\section{File formats, canonical form and test vectors}\label{sec:formats}


This appendix collects what an independent implementer needs in order to read the two certificate files
and to test a verifier written from \cref{alg:verify,alg:leaf}. It repeats a few definitions from the main
text so that it can be used on its own.

\subsection{Cells, lines, symmetries}\label{ssec:fmt-cells}

The cell $(x,y,z)\in\{0,\dots,5\}^3$ has the \emph{id} $36x+6y+z\in\{0,\dots,215\}$. All files use ids. A
\emph{line} is a set $\{(f_1(i),f_2(i),f_3(i)): i=0,\dots,5\}$ where each $f_j$ is a constant, $i$, or
$5-i$, and not all three are constant; there are $148$ lines ($108$ parallel to an axis, $36$ face diagonals,
$4$ space diagonals), and two cells lie on at most one common line (\cref{lem:geometry}). The $24$
\emph{busy} cells are those with $\min(x,5-x)=\min(y,5-y)=\min(z,5-z)$; their ids are
\begin{center}\small\ttfamily
0 5 30 35 43 46 61 64 86 87 92 93 122 123 128 129 151 154 169 172 180 185 210 215
\end{center}
A busy cell lies on $7$ lines; of the other cells, $144$ lie on $4$ lines and $48$ on $3$ lines.

The group $\G$ consists of the maps obtained by first permuting the three coordinates, then replacing
$x_j$ by $5-x_j$ for the $j$ in some subset of the coordinates, and then applying one bijection $s$ of
$\{0,\dots,5\}$ with $s(5-i)=5-s(i)$ to every coordinate. There are $6\cdot 8\cdot 48$ such descriptions and
$1{,}152$ distinct maps (\cref{lem:group}). $\G$ has three orbits on the cells, namely the three classes just
mentioned, with representatives $(0,0,0)$, $(0,0,1)$ and $(0,1,2)$.

\subsection{Sets, their order, and canonical forms}\label{ssec:fmt-canon}

A set of cells is written as the list of its ids in ascending order, separated by commas and without
spaces; the empty set is written as a single hyphen. Sets of the same size are ordered
lexicographically by these lists: $A<A'$ if, at the first position where the two ascending lists differ,
the id in the list of $A$ is the smaller one. (Equivalently: the smallest element of the symmetric
difference of $A$ and $A'$ belongs to $A$.)

The \emph{canonical form} of a pair $(B,F)$ of disjoint sets is the pair $(g(B),g(F))$, $g\in\G$, that is
smallest when the images of $B$ are compared first and, among the maps that minimize $g(B)$, the images
of $F$ second. A state $\Pst(B,F,t)$ or $\Qst(B,F,t)$ is canonical if $(B,F)$ is; $t$ and the type of the
state are not affected. Both files contain canonical states only. If $(B,F)$ is any state and $g$ is a map
with $(g(B),g(F))$ canonical, the table move $m$ listed for the canonical Q-state corresponds to the move
$g^{-1}(m)$ in the frame of $(B,F)$; any $g$ attaining the canonical form may be used (\cref{lem:canon}).

For a verifier it suffices to compute the canonical form by trying all $1{,}152$ maps. Test vectors (ids;
each line reads $B\mid F\ \to\ $ canonical $B\mid$ canonical $F$):
\begin{center}\small\ttfamily
\begin{tabular}{@{}l@{}}
43 | 0 \ $\to$\ 0 | 43\\
0,46,92,123,169,180,215 | - \ $\to$\ 0,5,61,93,122,154,210 | -\\
0,35,46,92,123,169,180,215 | - \ $\to$\ 0,5,61,93,122,154,210,215 | -\\
5,30,43,64,86,93,128,169,180 | 0,35,46,87,122,151,172,185,210\\
\qquad$\to$\ 0,35,43,64,86,93,123,154,210 | 5,30,61,122,129,151,172,180,215
\end{tabular}
\end{center}
The second and third are the sets $\Tseven$ and $\Teight$ of \cref{ex:T7,ex:T8}; the last is the final
position of Line~B in \cref{ssec:walk}.

\subsection{The strategy table}\label{ssec:fmt-table}

The file \fn{results/final/strategy.tbl} is a text file with $93{,}164$ lines, one for each Q-state of the
certificate, in no particular order and without a header. A line has four fields separated by single
spaces:
\begin{center}
\texttt{t\ \ B\ \ F\ \ m}
\end{center}
where $t$ is the round in decimal, $B$ and $F$ are sets in the notation above, and $m$ is a cell id. The
line says: \emph{in the canonical Q-state $\Qst(B,F,t)$ --- Maker has made his $t$-th move, Breaker owns the
$t-1$ stones $B$, the Maker stones $F$ are revealed, $|F|\le t$ --- Breaker's table move is the cell $m$}.
The move is given in the frame of the canonical state and satisfies $m\notin B\cup F$. No key $(t,B,F)$
occurs twice. Three lines of the file (the first, line $2{,}354$, and line $9{,}791$):
\begin{verbatim}
1 - - 0
5 0,46,169,215 - 92
7 0,35,46,92,123,169 5,30,43,180,185,210,215 50
\end{verbatim}
In coordinates: the first is Breaker's opening move $(0,0,0)$ in the state $\Qst(\emptyset,\emptyset,1)$.
The second says that with $B=\{(0,0,0),(1,1,4),(4,4,1),(5,5,5)\}$ and nothing revealed, Breaker plays
$(2,3,2)$ in round $5$. In the third, Breaker owns $(0,0,0)$, $(0,5,5)$, $(1,1,4)$, $(2,3,2)$, $(3,2,3)$,
$(4,4,1)$, the seven revealed Maker stones are $(0,0,5)$, $(0,5,0)$, $(1,1,1)$, $(5,0,0)$, $(5,0,5)$,
$(5,5,0)$, $(5,5,5)$, and Breaker's move in round $7$ is the non-busy cell $(1,2,2)$.

\subsection{The leaf file}\label{ssec:fmt-leaves}

The file \fn{results/final/leaves.txt} has $106{,}723$ lines with three fields,
\begin{center}
\texttt{t\ \ B\ \ F}
\end{center}
meaning: \emph{the canonical P-state $\Pst(B,F,t)$ --- $t$ rounds are complete, Breaker owns the $t$ stones
$B$, the Maker stones $F$ are revealed --- is a leaf}, i.e.\ $\LEAF(B,F,t)$ is claimed to hold and the
state is not expanded. Three lines (the first, line $8{,}185$, and the last):
\begin{verbatim}
9 0,1,5,30,43,93,122,172,210 7,13,14,19,21,25,28,31,35
8 0,5,61,93,122,154,210,215 -
9 0,35,43,64,86,93,123,154,210 5,30,61,122,129,151,172,180,215
\end{verbatim}
In the first, Breaker owns $(0,0,0)$, $(0,0,1)$, $(0,0,5)$, $(0,5,0)$, $(1,1,1)$, $(2,3,3)$, $(3,2,2)$,
$(4,4,4)$, $(5,5,0)$, and $F$ consists of the nine cells $(0,1,1)$, $(0,2,1)$, $(0,2,2)$, $(0,3,1)$,
$(0,3,3)$, $(0,4,1)$, $(0,4,4)$, $(0,5,1)$, $(0,5,5)$: Maker has made two threats, on the line $\{(0,i,1)\}$ and on the diagonal
$\{(0,i,i)\}$ of the face $x=0$, which share the stone $(0,1,1)$; Breaker's forced replies are $(0,0,1)$ and
$(0,0,0)$, and all nine Maker stones are revealed ($\theta=0$). The second line is the canonical form of
$\Pst(\Teight,\emptyset,8)$, the only leaf with $\theta=8$. The third is the final state of Line~B of
\cref{ssec:walk}.

A verifier must not assume that the files are correct; \cref{alg:verify} only uses them as hints. It is
reasonable to check, in addition, that every line is well formed ($|B|=t-1$ resp.\ $|B|=t$, $|F|\le t$,
$B\cap F=\emptyset$, $m\notin B\cup F$, ascending lists), canonical, and not repeated; this holds for both
files.

\subsection{Numbers a verifier should reproduce}\label{ssec:fmt-targets}

A correct implementation of \cref{alg:verify} with $\Rmax=10$ accepts the two files and visits
$246{,}299$ states: $153{,}135$ P-states, of which $46{,}412$ are expanded and $106{,}723$ are leaves, and
$93{,}164$ Q-states. Per round $t=0,\dots,10$ (\cref{tab:tree}):
\begin{center}\footnotesize\setlength{\tabcolsep}{3.2pt}
\begin{tabular}{@{}lrrrrrrrrrrr@{}}
\toprule
$t$ & 0 & 1 & 2 & 3 & 4 & 5 & 6 & 7 & 8 & 9 & 10\\
\midrule
P-states & 1 & 2 & 5 & 14 & 35 & 181 & 1{,}718 & 7{,}227 & 34{,}801 & 93{,}124 & 16{,}027\\
leaves & 0 & 0 & 0 & 0 & 0 & 0 & 0 & 0 & 3{,}382 & 87{,}314 & 16{,}027\\
Q-states & 0 & 2 & 5 & 14 & 35 & 101 & 446 & 3{,}487 & 14{,}362 & 62{,}381 & 12{,}331\\
threat children (c) & 0 & 0 & 0 & 0 & 0 & 824 & 4{,}075 & 14{,}460 & 47{,}020 & 72{,}352 & 4{,}767\\
danger children (d) & 0 & 0 & 0 & 0 & 0 & 0 & 0 & 0 & 20{,}416 & 6{,}060 & 1{,}799\\
\VOID{} children (ii) & 0 & 0 & 0 & 0 & 0 & 0 & 1 & 6 & 54 & 88 & 2\\
\bottomrule
\end{tabular}
\end{center}
The conventions are those of \cref{tab:tree}: states are counted as canonical forms; a threat or danger
child is counted, at the round of the child, for every pair of an expanded canonical P-state and one of
its non-\VOID{} children (the children of one state are pairwise distinct, but they may be symmetric to
each other or to children of other states); the last row counts the \VOID{} children of type (ii) at the
round of their parent Q-state. In total there are $143{,}498$ threat children, $28{,}275$ danger children
and $54{,}482$ children of type (ii), of which $151$ are \VOID; no child of type (c) or (d) is \VOID. Every
line of both files is reached. The leaves split by $\theta=t-|F|=0,\dots,8$ into
$64{,}841$, $32{,}703$, $7{,}267$, $1{,}434$, $343$, $105$, $25$, $4$, $1$ (\cref{tab:leaf-theta}), and $434$ table moves are
non-busy cells ($16$, $399$, $11$, $8$ in rounds $7$ to $10$).

\subsection{Test vectors for the leaf test}\label{ssec:fmt-leaf}

Let
\[
\Tseven=\{(0,0,0),(5,5,5),(1,1,4),(4,4,1),(2,3,2),(3,2,3),(5,0,0)\},\qquad
\Teight=\Tseven\cup\{(0,5,5)\},
\]
with ids \texttt{0,46,92,123,169,180,215} and \texttt{0,35,46,92,123,169,180,215}. \Cref{tab:leafvectors}
lists values of $\LEAF(B,F,k)$ that an implementation of \cref{alg:leaf} should reproduce; an instance is
written as ``$k$ $B$ $F$'' in the notation of the files. Each value was obtained with three independent
implementations (the two verifiers and the library of the search), which also agree on which of the
three conditions (L1)--(L3) of \cref{def:leaf} fails. The last column gives, for comparison, the number
$\worst(B,F,k)$ of \cref{sec:ip} computed with the integer program; the pairing certificate
$\CC(B,F,k)$ holds if and only if it is $0$.

\begin{table}[ht]
\centering\footnotesize\setlength{\tabcolsep}{2.5pt}
\begin{tabular}{@{}llcl@{}}
\toprule
instance ``$k$ $B$ $F$'' (ids) & description & $\LEAF$ & $\worst$\\
\midrule
\texttt{7 0,46,92,123,169,180,215 -} & $\Tseven$, $k=7$ & false (3) & 6\\
\texttt{8 0,35,46,92,123,169,180,215 -} & $\Teight$, $k=8$ & true & 0\\
\texttt{9 0,35,46,92,123,169,180,215 -} & $\Teight$, $k=9$ & false (3) & 1\\
\texttt{8 0,35,46,92,123,169,180,215 }$h$ & $\Teight$, $F=\{h\}$, each of the $16$ busy $h\notin\Teight$ & true & 0\\
\texttt{8 0,35,43,46,92,123,169,180 215} & $\Teight-(5,5,5)+(1,1,1)$, $F=\{(5,5,5)\}$ & true & 0\\
\texttt{8 0,5,35,43,46,61,86,92} & antipode stealing after $8$ rounds & false (2) & 6\\
\quad\texttt{123,129,169,172,180,210,215} & & &\\
\texttt{9 0,5,35,43,46,61,86,87,92} & antipode stealing after $9$ rounds & true & 0\\
\quad\texttt{123,129,154,169,172,180,210,215} & & &\\
\texttt{10 0,5,30,35,43,46,61,86,87,92} & antipode stealing after $10$ rounds & true & 0\\
\quad\texttt{123,128,129,154,169,172,180,210,215} & & &\\
\texttt{9 0,8,35,46,92,123,169,180,215} & threat: $\Teight\cup\{(0,1,2)\}$, $F=\{(x,1,2)\colon x\ge1\}$ & true & 0\\
\quad\texttt{44,80,116,152,188} & & &\\
\texttt{8 0,8,46,92,123,169,180,215} & threat: $\Tseven\cup\{(0,1,2)\}$, the same $F$ & false (3) & 2\\
\quad\texttt{44,80,116,152,188} & & &\\
\texttt{0 - -} & the empty board & false (2) & 80\\
\bottomrule
\end{tabular}
\caption{Test vectors for $\LEAF(B,F,k)$. ``false ($i$)'' means that condition (L$i$) of \cref{def:leaf}
is the first to fail: (L2) no pairing for $(B,F)$; (L3) an $F$-connected set of at least three live lines
has surplus less than $k-|F|$. The $16$ cells $h$ in the fourth row have the ids
5, 30, 43, 61, 64, 86, 87, 93, 122, 128, 129, 151, 154, 172, 185, 210.}
\label{tab:leafvectors}
\end{table}

Some explanations (see also \cref{fact:T8-leaf}). For $\Tseven$ the set of all $104$ live lines is
$\emptyset$-connected and has surplus $216-7-208=1=\slack(\Tseven)<7$. For $\Teight$ the $100$ live lines
have surplus $216-8-200=8$, which is the minimum over all $\emptyset$-connected sets of at least three
live lines; hence $\LEAF(\Teight,\emptyset,8)$ holds and $\LEAF(\Teight,\emptyset,9)$ fails. In the ``antipode stealing''
instances Breaker has played $(0,0,0)$, $(1,1,4)$, $(2,3,2)$, $(0,5,5)$, $(1,1,1)$, $(2,2,2)$, $(0,0,5)$,
$(1,4,1)$ and then $(2,2,3)$ and $(0,5,0)$, and $F$ consists of the antipodes of all his stones except the
last one. In the two threat instances Maker holds five cells of the line $\{(x,1,2)\}$ and Breaker has been
forced to $(0,1,2)$. On the empty board $148$ lines would need $296$ cells, and the deficiency
$296-216=80$ is the value of $\worst$.

Further values can be generated with the programs of the artifact: \fn{bin/phicheck} and
\fn{bin/treecheck} with the option \fn{--leafcheck} read instances ``$k$ $B$ $F$'' from standard input and
print the value of $\LEAF$, the failing condition and a witness.

\subsection{Orbit counts and the leaf table}\label{ssec:fmt-orbits}

The numbers of orbits of $k$-subsets of the $24$ busy cells under $\G$ are
\[
527,\quad 1{,}057,\quad 1{,}669,\quad 2{,}468\qquad\text{for } k=7,\,8,\,9,\,10 .
\]
All four were computed by Burnside's lemma and by enumeration of the canonical representatives
(\fn{bin/test_canon}), and the first three also by a second program, in Python. The files
\fn{results/leaftable/phi_k8.jsonl}, \fn{phi_k9.jsonl} and \fn{phi_k10.jsonl} contain one JSON record per
orbit ($1{,}057$, $1{,}669$ and $2{,}468$ lines) with the canonical representative $B$, its slack, and the value of
$\LEAF(B,\emptyset,k)$. The counts are those of \cref{tab:leaftable}:
$6$, $947$, $2{,}185$ orbits have $\slack(B)\ge k$, and $6$, $172$, $0$ satisfy $\LEAF(B,\emptyset,k)$, for
$k=8,9,10$. The six representatives for $k=8$ are listed in \cref{ssec:leaftable}; in ids they are
\begin{center}\small\ttfamily
\begin{tabular}{@{}ll@{}}
0,5,43,61,93,122,154,210 & 0,5,43,92,93,122,172,210\\
0,5,43,61,93,122,154,215 & 0,5,43,92,93,122,172,215\\
0,5,43,61,93,122,172,215 & 0,5,61,93,122,154,210,215
\end{tabular}
\end{center}

\subsection{The minimum number of stones admitting a pairing}\label{ssec:fmt-minstones}

A last set of regression values concerns the $n\times n\times n$ board with its $((n+2)^3-n^3)/2$ lines:
what is the smallest number of Breaker stones $B$ such that a pairing for $(B,\emptyset)$ exists? The
program \fn{src/minstones.cpp} gives
\begin{center}\small
\begin{tabular}{@{}lrrrr@{}}
\toprule
$n$ & 5 & 6 & 7 & 8\\
\midrule
lines & 109 & 148 & 193 & 244\\
minimum number of stones & 9 & 7 & 3 & 0\\
\bottomrule
\end{tabular}
\end{center}
(\fn{results/regression/minstones.json}). What is established there is the following. \emph{Lower
bounds:} a pairing needs two cells for every live line, so $2|\live{B}|\le n^3-|B|$; for every size $j$
below the stated minimum the program searches all $j$-subsets (up to the symmetry of the first stone,
and with a bound on the number of lines that the remaining stones can kill) and finds that none satisfies
even this counting condition. For $n=6$, six stones kill at most $39$ lines, while $43$ would have to be
killed. \emph{Upper bounds:} for each $n$ the program exhibits a set of the stated size together with an
explicit pairing, which is checked against the definition by a separate routine
(\fn{results/regression/minstones_pairings.json}); for $n=6$ the set is
$\{(0,0,0),(0,0,5),(1,4,1),(2,3,3),(3,2,2),(4,1,4),(5,5,0)\}$, the canonical form of $\Tseven$, and for $n=8$
the empty set: the $8\times8\times8$ board admits a pairing of its $244$ lines without any stone ($488$ of
the $512$ cells are used).
